\scanpage{95}
I) $\alpha\ne\pm2$ \quad (valeurs propres distinctes de $u_V$, \unsure{exste|existe} centre)

\deleted{Sur l'extension \unsure{$k_\zeta$|finie} \unsure{extension} $k'$ de $k$}
Si $k$ contient les racines $\unsure{\zeta|\rho},\zeta'=\zeta^{-1}$ de l'équ aux valeurs propres, alors \unsure{can.|cas|on est} isom. à \deleted{$\Pi_\zeta$} $\varphi_{\unsure{\zeta}}$ \dots{} Ordre de $u_V$ est égal à celui de $u_E$ \unsure{de même|du moins|ou encore} à celui de $\zeta$.
\deleted{II) $\alpha=-2$, $\alpha\ne2$}
\deleted{II) Cas $\alpha=\unsure{\pm|-}1$ i.e. valeurs propres confondues} \deleted{\unsure{II)} Cas $\alpha={}$} \deleted{II) Cas $\alpha=-2$, $\alpha\ne2$ il y a}
\[
\deleted{u=\begin{pmatrix}\unsure{1{-}2|-2}&-1&1\\1&0&0\\0&0&1\end{pmatrix}}
\]
\[
\begin{aligned}
\text{II)}\quad\alpha&=-2\qquad \zeta=\zeta'=-1&
\deleted{u_V^{\unsure{2}}=\mathrm{id}}\quad u_V^2&=\text{transvection}\\
&&&\text{d'ordre }p\\
&&u_V&\ \text{d'ordre }2p\\
&&u_E&\ \text{d'ordre }2p\text{ si }p\ne2\\
&&&(\text{car }\exists\text{ centre})
\end{aligned}
\]
$\overline{\mathrm M}$ conclusion si $p=2$. En effet, on a alors $\alpha=-2=0$, or le cas $\alpha=0$ est celui du \uline{carré} ($\zeta^2=-1$ \unsure{§|!}\dots)
\[
\begin{aligned}
\deleted{\sigma_{\unsure{0}}(x,y}\quad u(x,y)&=(1-y,x)\\
u^2(x,y)&=(1-x,1-y)\\
u^4(x,y)&=(x,y)\qquad\cdots
\end{aligned}
\]
Si $p\ne2$, \inserted{i.e. $\alpha\ne2$,} on trouve, en posant
\[
\left\{\begin{aligned}X&=x\unsure{-|+}\frac{1}{2-\alpha}\\Y&=y\unsure{-|+}\frac{1}{\unsure{2|7}-\alpha}\end{aligned}\right.
\]
\[
\begin{aligned}
\left\{\begin{aligned}\sigma_0(x,y)&=(-X+\alpha Y,Y)\\\sigma_1(X,Y)&=(Y,X)\\u(X,Y)&=(-Y+\alpha X,X)=-(\mathrm{id}+a'\otimes a)(X,Y)\end{aligned}\right.\\
a'(X,Y)=X+Y\qquad a=(1,-1)
\end{aligned}
\]
% archive sheet number: 94
